\chapterphoto{UCSC McHenry 图书馆与 Media Theater 之间的小径}{photo-chap2}
\chapter{PID 控制器}

PID 控制器是一种常用的反馈控制器，由比例（proportional）、积分（integral）、微分（derivative）三项组成，PID 因此得名。本章将逐项构建 PID 控制器的定义，同时尽力为每一项的行为建立直觉。

首先，我们约定一些 PID 控制器的术语。参考值在这里称为设定值（setpoint，即期望位置），输出称为过程变量（process variable，即测量到的位置）。下面是相关量的一些常见变量命名约定：

\begin{center}
\begin{tabular}{ll}
  $r(t)$ & 设定值 \\
  $u(t)$ & 控制输入 \\
  $e(t)$ & 误差 \\
  $y(t)$ & 输出 \\
\end{tabular}
\end{center}

其中误差 $e(t) = r(t) - y(t)$。

对于已经熟悉 PID 控制的读者需要说明的是：本书的解读与经典的``过去、现在、未来''误差直觉并不一致。我们将从现代控制理论的视角出发，把比例控制器应用于我们所关心的不同物理量。这样能够更完整地解释导数项在设定值恒定和变化时的行为，而且这种直觉也会自然延续到本书后面介绍的现代控制方法中。

概括地说：比例项将\textbf{位置}误差驱动到零，微分项将\textbf{速度}误差驱动到零，而积分项随时间累积设定值曲线与输出曲线之间的面积（即位置误差的积分），并把当前的累积值加到控制输入上。下面我们对每一项展开详细讨论。

\section{比例项}
\label{sec:ptermm}

比例项将位置误差驱动到零。

\begin{definition}[比例控制器]
\begin{equation}
  u(t) = K_p e(t)
\end{equation}
其中 $K_p$ 是比例增益，$e(t)$ 是当前时刻 $t$ 的误差。
\end{definition}

图~\ref{fig:pblock} 展示了由 P 控制器控制的系统的框图。

\begin{figure}[htbp]
  \centering
  \includegraphics[width=0.85\textwidth]{fig2-1.png}
  \caption{P 控制器框图}
  \label{fig:pblock}
\end{figure}

比例增益的行为就像``用软件定义的弹簧''，把系统拉向期望的位置。回想物理课上学过的弹簧模型 $F = -kx$，其中 $F$ 是施加的力，$k$ 是比例常数，$x$ 是离开平衡点的位移。这个式子也可以写成 $F = k(0 - x)$，其中 $0$ 是平衡点。如果让平衡点成为反馈控制器的设定值，那么两个方程便一一对应：

\begin{align*}
  F &= k(r - x) \\
  u(t) &= K_p e(t) = K_p \left( r(t) - y(t) \right)
\end{align*}

也就是说，比例控制器把系统输出拉向设定值的``力''与误差成正比——正如弹簧一样。

\section{微分项}
\label{sec:dterm}

微分项将速度误差驱动到零。

\begin{definition}[PD 控制器]
\begin{equation}
  u(t) = K_p e(t) + K_d \frac{de}{dt}
\end{equation}
其中 $K_p$ 是比例增益，$K_d$ 是微分增益，$e(t)$ 是当前时刻 $t$ 的误差。
\end{definition}

图~\ref{fig:pdblock} 展示了由 PD 控制器控制的系统的框图。

\begin{figure}[htbp]
  \centering
  \includegraphics[width=0.9\textwidth]{fig2-2.png}
  \caption{PD 控制器框图}
  \label{fig:pdblock}
\end{figure}

PD 控制器包含一个位置比例控制器（$K_p$）和一个速度比例控制器（$K_d$）。速度设定值是由位置设定值随时间的变化隐式给出的。图~\ref{fig:pdelevator} 展示了电梯机构上的一个例子。

为了证明 PD 控制器其实就是两个比例控制器，我们对 PD 控制器的方程进行变形。设
\[
  u_k = K_p e_k + K_d \frac{e_k - e_{k-1}}{\Delta t}
\]
其中 $u_k$ 是第 $k$ 个时间步的控制输入，$e_k$ 是第 $k$ 个时间步的误差，$\Delta t$ 是时间步长。$e_k$ 定义为 $e_k = r_k - y_k$，其中 $r_k$ 是第 $k$ 个时间步的设定值，$y_k$ 是第 $k$ 个时间步的输出。

\begin{align*}
  u_k &= K_p (r_k - y_k) + K_d \frac{(r_k - y_k) - (r_{k-1} - y_{k-1})}{\Delta t} \\
  u_k &= K_p (r_k - y_k) + K_d \frac{r_k - y_k - r_{k-1} + y_{k-1}}{\Delta t} \\
  u_k &= K_p (r_k - y_k) + K_d \frac{r_k - r_{k-1} - y_k + y_{k-1}}{\Delta t} \\
  u_k &= K_p (r_k - y_k) + K_d \frac{(r_k - r_{k-1}) - (y_k - y_{k-1})}{\Delta t} \\
  u_k &= K_p (r_k - y_k) + K_d \left( \frac{r_k - r_{k-1}}{\Delta t} - \frac{y_k - y_{k-1}}{\Delta t} \right)
\end{align*}

注意，$\frac{r_k - r_{k-1}}{\Delta t}$ 是设定值的速度，$\frac{y_k - y_{k-1}}{\Delta t}$ 是系统速度的估计值。这意味着 PD 控制器的 $K_d$ 项把估计速度驱动到设定值速度。

如果设定值是恒定的，隐含的速度设定值就是零，因此当系统在运动时，$K_d$ 项会使它减速。它的作用就像``用软件定义的阻尼器''。这类阻尼器常见于闭门器上，其阻尼力随速度线性增大。

\begin{figure}[htbp]
  \centering
  \includegraphics[width=0.62\textwidth]{fig2-3.png}
  \caption{电梯机构上的 PD 控制器}
  \label{fig:pdelevator}
\end{figure}

\section{积分项}
\label{sec:iterm}

积分项随时间累积设定值曲线与输出曲线之间的面积（即位置误差的积分），并把当前的累积值加到控制输入上。累积两条曲线之间的面积称为积分。

\begin{definition}[PI 控制器]
\begin{equation}
  u(t) = K_p e(t) + K_i \int_0^t e(\tau) \, d\tau
\end{equation}
其中 $K_p$ 是比例增益，$K_i$ 是积分增益，$e(t)$ 是当前时刻 $t$ 的误差，$\tau$ 是积分变量。
\end{definition}

积分从时刻 $0$ 一直积到当前时刻 $t$。我们用 $\tau$ 表示积分变量，因为在积分过程中需要一个能够取多个值的变量，而 $t$ 已经被定义为当前时间，不能再使用。

图~\ref{fig:piblock} 展示了由 PI 控制器控制的系统的框图。

\begin{figure}[htbp]
  \centering
  \includegraphics[width=0.9\textwidth]{fig2-4.png}
  \caption{PI 控制器框图}
  \label{fig:piblock}
\end{figure}

当系统在稳态下已经非常接近设定值时，比例项可能太小，不足以把输出一路拉到设定值，而此时微分项为零。这就可能造成稳态误差，如图~\ref{fig:sse} 所示。

消除稳态误差的一种常用方法是对误差进行积分，并将其加到控制输入上。这会使控制量不断增大，直到系统收敛为止。图~\ref{fig:sse} 展示了飞轮出现稳态误差的例子；图~\ref{fig:pi-nosse} 展示了在飞轮控制器中加入积分器之后如何消除了稳态误差。然而，积分增益过大会导致超调，如图~\ref{fig:pi-overshoot} 所示。

\begin{tip}
消除稳态误差还有更好的办法，例如使用前馈，或者利用对系统的其他知识来约束积分控制作用的时机。我们讲到现代控制理论时会更详细地讨论这些内容。
\end{tip}

\begin{figure}[htbp]
  \centering
  \includegraphics[width=0.62\textwidth]{fig2-5.png}
  \caption{飞轮上带有稳态误差的 P 控制器}
  \label{fig:sse}
\end{figure}

\begin{figure}[htbp]
  \begin{minipage}[t]{0.48\textwidth}
    \centering
    \includegraphics[width=\linewidth]{fig2-6.png}
    \caption{飞轮上无稳态误差的 PI 控制器}
    \label{fig:pi-nosse}
  \end{minipage}\hfill
  \begin{minipage}[t]{0.48\textwidth}
    \centering
    \includegraphics[width=\linewidth]{fig2-7.png}
    \caption{$K_i$ 增益过大导致超调的飞轮 PI 控制器}
    \label{fig:pi-overshoot}
  \end{minipage}
\end{figure}

\section{PID 控制器的定义}
\label{sec:piddef}

把这三项组合起来，就得到了 PID 控制器的典型定义。

\begin{definition}[PID 控制器]
\begin{equation}
  u(t) = K_p e(t) + K_i \int_0^t e(\tau) \, d\tau + K_d \frac{de}{dt}
\end{equation}
其中 $K_p$ 是比例增益，$K_i$ 是积分增益，$K_d$ 是微分增益，$e(t)$ 是当前时刻 $t$ 的误差，$\tau$ 是积分变量。
\end{definition}

图~\ref{fig:pidblock} 展示了由 PID 控制器控制的系统的框图。

\begin{figure}[htbp]
  \centering
  \includegraphics[width=0.9\textwidth]{fig2-8.png}
  \caption{PID 控制器框图}
  \label{fig:pidblock}
\end{figure}

\section{响应类型}
\label{sec:responsetypes}

由 PID 控制器驱动的系统一般有三种响应类型：欠阻尼、过阻尼和临界阻尼。如图~\ref{fig:responses} 所示。

\begin{figure}[htbp]
  \centering
  \includegraphics[width=0.62\textwidth]{fig2-9.png}
  \caption{PID 控制器的响应类型}
  \label{fig:responses}
\end{figure}

对于图~\ref{fig:responses} 中的阶跃响应：上升时间（rise time）是系统在施加阶跃输入后首次到达参考值所花的时间；调节时间（settling time）是系统在施加阶跃输入后在参考值处稳定下来所花的时间。

欠阻尼响应在稳定之前会围绕参考值振荡。过阻尼响应上升缓慢，且不会越过参考值。临界阻尼响应在不围绕参考值振荡（即先过冲后回落）的前提下，具有最短的上升时间。

\section{手动整定}
\label{sec:manualtuning}

以下步骤适用于位置式 PID 控制器。速度式 PID 控制器通常不需要 $K_d$。

\begin{enumerate}
  \item 将 $K_p$、$K_i$ 和 $K_d$ 置零。
  \item 增大 $K_p$，直到输出开始在设定值附近振荡。
  \item 在不引起系统响应抖动的前提下，尽可能增大 $K_d$。
\end{enumerate}

如果设定值遵循梯形曲线（见第 15.1 节），整定会容易得多。将位置设定值、速度设定值、测量位置和测量速度都绘制出来。速度设定值可以通过对位置设定值做数值微分得到（即 $v_{\text{desired},k} = \frac{r_k - r_{k-1}}{\Delta t}$）。先增大 $K_p$，直到位置跟踪良好；然后增大 $K_d$，直到速度跟踪良好。

如果控制器最终稳定在高于或低于设定值的输出上，可以增大 $K_i$，使控制器在合理的时间内到达设定值。

\begin{tip}
如果有模型可用（基于输出的控制并非如此），应优先使用前馈（第 7.8 节）或在被控对象中引入积分器（第 6.7 节），而不是在控制器中使用积分器（例如 PID 的 I 项）。积分器只应用于补偿无法建模的动态。
\end{tip}

注意，如果 $K_i$ 过大，就可能发生积分饱和（integral windup）：设定值发生大幅变化后，积分项累积的误差可能超过最大控制输入。结果，系统会超调并持续增大，直到这部分累积的误差被释放为止。

\section{局限性}
\label{sec:limitations}

当对系统动态没有任何先验知识时，PID 的启发式整定方法是合理的选择。然而，如果已知系统的动力学模型，就可以设计出响应好得多的控制器。此外，PID 只适用于单输入单输出（SISO）系统；我们将在本书第二部分介绍多输入多输出（MIMO）系统的控制方法。
